\weeknum{5}
\topic[Isostasy]{Isostasy: Ocean--Continent Structure \& Isostatic Gravity}

\workshopstart

\vspace{-0.75em}
\noindent\textbf{Duration:} 3 hours \hfill \textbf{Setting:} Workshop room (small-group whiteboard work)

\section*{Preparation}

\begin{description}
  \item[Pre-reading:] Course Notes, Ch. 7
  \item[Lecture videos:]
  \item[Notes:]
\end{description}

\section*{Purpose}

This workshop develops a quantitative understanding of isostatic equilibrium and its gravity expression by combining column-based reasoning, ocean--continent structure, and slab gravity approximations. Students work from a single general mass-balance principle throughout, applying it first to a crust-only ocean--continent comparison, then to a layered lithosphere, then to gravity. Airy and Pratt are introduced as named patterns within this balance, not as separate starting formulas.

\section*{Learning Outcomes}

By the end of this workshop, you should be able to:

\begin{itemize}
    \item Explain isostatic equilibrium as equal lithostatic pressure (or mass per unit area) at a compensation depth, for columns of arbitrary structure.
    \item Diagnose when a model result (e.g.\ an unrealistic density) indicates a missing physical component, and propose what is missing.
    \item Recognize whether a given compensation pattern is consistent with the Airy or Pratt end-member, and explain why.
    \item Derive the isostatic gravity correction from mass balance and the infinite slab formula, and explain why large topography does not necessarily produce large gravity anomalies.
    \item Explain compensation as a spectrum rather than a binary state, and use misfit between predicted and observed gravity to reason about the degree of compensation.
\end{itemize}

\section*{Session Flow (\timelinelength\ min)}

\timeline[slot]{Warm-Up: Icebergs}{10}

Isostasy is a state of equilibrium related to the buoyancy of the lithosphere over a geophysically fluid mantle. In this exercise, we'll examine the buoyancy of icebergs and discuss the similarity to the lithosphere that we'll be studying.

\timeline[slot]{Mini-Lecture: The General Isostatic Balance}{10}

\begin{itemize}
    \item Isostatic equilibrium as pressure balance at a common compensation depth
    \item Handling columns of unequal height
\end{itemize}

The full statement of this balance, and how to apply it to columns of unequal height, is given in the theory box at the start of the next activity.

\timeline[item]{Oceanic vs.\ Continental Crust, Part 1}{30}

Each group balances an oceanic and a continental column, assuming compensation occurs solely through crustal thickness, and solves for the oceanic crust density required for equilibrium. See the activity sheet for the theory and setup.

\timeline[slot]{Mini-Lecture: End-Member Terminology}{10}

\begin{itemize}
    \item Naming the two patterns you just used: Airy (uniform density, variable thickness) and Pratt (uniform thickness, variable density)
    \item Schematic comparison of the two patterns, and why real structure usually mixes both
\end{itemize}

\timeline[item]{Oceanic vs.\ Continental Crust, Part 2}{25}

The model is extended to include a lithospheric mantle beneath both columns. Groups rebalance the columns, estimate the density contrast between oceanic and continental lithospheric mantle, and relate it to mineralogy (e.g.\ garnet content). See the activity sheet.

\timeline[item]{Isostatic Gravity Correction}{40}

Working from the same balance, groups derive the compensating root thickness, apply the infinite slab formula (Week~4) to topography and root, and show that their gravity effects cancel under full compensation. See the activity sheet for the full derivation sequence.

\timeline[slot]{Numerical Demonstration}{30}

The instructor builds a lithospheric cross section in \texttt{isostasy\_gui.py} and computes its gravity response under varying degrees of compensation and layered structure. Groups work through the accompanying activity sheet as the demonstration proceeds.

\timeline[slot]{Wrap-up}{15}

Each group prepares a brief summary addressing:
\begin{itemize}
    \item One case in this workshop where an unrealistic result told you something was missing from the model, and what that turned out to be.
    \item Whether a compensation pattern you examined today was closer to Airy, Pratt, or a mix, and why.
    \item One reason a compensated load can produce little or no gravity anomaly, and one reason a real gravity residual might not be zero.
    \item One observation (e.g.\ seismology or petrology) that could help distinguish between isostatic models that fit the same gravity data equally well.
\end{itemize}

\timeline[slot]{Preview: Magnetics}{10}

Next week we move from gravity to magnetics. Like gravity, magnetic methods respond to a physical property contrast---in this case, magnetic susceptibility rather than density. However, magnetics differs in important ways:
\begin{itemize}
    \item magnetic susceptibility is controlled by trace rather than bulk mineralogy;
    \item rock magnetization is a vector (not a scalar);
    \item the response depends on the direction of the regional field; and
    \item the source contribution changes sign with polarity.
\end{itemize}
We will use hands-on demonstrations with magnets and rock samples to build intuition before examining how crustal structure is recorded in magnetic anomalies.

\workshopend
\notepage
\notepage

\subimport{activities/}{activity_ocean_continent}
\subimport{activities/}{activity_isostatic_correction}
\subimport{activities/}{activity_isostatic_demo}